\DocumentMetadata{
  pdfversion=2.0,
  pdfstandard=ua-2,
  lang=en-US,
  tagging=on,
  tagging-setup={math/setup=mathml-SE,tabsorder=structure}
}
\documentclass[11pt,letterpaper]{article}
\usepackage[margin=0.68in,includefoot]{geometry}
\usepackage{newcomputermodern}
\usepackage{amsmath,amscd,mathtools}
\providecommand{\square}{\mdlgwhtsquare}
\usepackage[hidelinks]{hyperref}
\hypersetup{
  pdftitle={Algebra PhD Exam, August 2012},
  pdfauthor={University of Florida Department of Mathematics},
  pdfsubject={Graduate mathematics examination},
  pdfdisplaydoctitle=true,
  bookmarksopen=true
}
\setcounter{secnumdepth}{0}
\setlength{\parindent}{0pt}
\setlength{\parskip}{0.28em}
\setlength{\emergencystretch}{3em}
\setlength{\leftmargini}{2.15em}
\setlength{\leftmarginii}{2.65em}
\setlength{\leftmarginiii}{3em}
\renewcommand{\labelenumi}{\textbf{\theenumi.}}
\pagestyle{empty}
\raggedbottom
\allowdisplaybreaks
\newenvironment{examproblems}[1]{%
  \begin{enumerate}%
  \setcounter{enumi}{#1}%
  \setlength{\topsep}{0.35em}%
  \setlength{\partopsep}{0pt}%
  \setlength{\itemsep}{0.48em}%
  \setlength{\parsep}{0.18em}%
}{\end{enumerate}}
\newenvironment{examsubparts}{%
  \begin{enumerate}%
  \setlength{\topsep}{0.18em}%
  \setlength{\partopsep}{0pt}%
  \setlength{\itemsep}{0.16em}%
  \setlength{\parsep}{0.1em}%
}{\end{enumerate}}
\newenvironment{examinstructions}{%
  \par\begingroup\itshape
}{%
  \par\endgroup\vspace{0.18em}
}
\makeatletter
\renewcommand\section{\@startsection{section}{1}{0pt}%
  {-1.25ex plus -.25ex minus -.1ex}%
  {0.55ex plus .1ex}%
  {\normalfont\large\bfseries}}
\renewcommand{\@maketitle}{%
  \newpage\null\vskip -1.2em%
  {\footnotesize\bfseries
    \href{https://www.ufl.edu/}{University of Florida}, \href{https://math.ufl.edu/}{Department of Mathematics}\par}%
  \vskip 0.18em%
  \tagstructbegin{tag=Title}%
    {\LARGE\bfseries\href{https://gma.math.ufl.edu/exams/algebra-phd/}{Algebra PhD Exam}, August 2012\par}%
  \tagstructend%
  \vskip 0.35em%
  \tagmcbegin{artifact}%
    \hrule height 1pt%
  \tagmcend%
  \vskip 0.55em%
}
\makeatother
\title{Algebra PhD Exam, August 2012}
\author{}
\date{}
\begin{document}
\maketitle
\thispagestyle{empty}
\begin{examinstructions}
Answer seven problems. (If you turn in more, the first seven will be graded.) Within reason, you may quote theorems as long as you state them clearly.
\end{examinstructions}
\begin{examproblems}{0}
\item
This problem has two parts.
\begin{examsubparts}
\item[\textnormal{(a)}]
Show that for all \(n \geq 0\) there exists an irreducible polynomial \(f \in \mathbb{Q}[x]\) of degree greater than~\(n\) that does not satisfy the hypotheses of Eisenstein’s criterion.
\item[\textnormal{(b)}]
Give an example, with proofs, of an irreducible polynomial which is not separable.
\end{examsubparts}
\item
Let \(n \geq 1\) and let~\(F\) be a splitting field of \(x^n-1 \in \mathbb{Q}[x]\).
\begin{examsubparts}
\item[\textnormal{(a)}]
Prove that the Galois group of~\(F\) over~\(\mathbb{Q}\) is abelian.
\item[\textnormal{(b)}]
When \(n=8\) show that the Galois group is of order~\(4\) and not cyclic.
\end{examsubparts}
\item
Prove that the Galois group of \(x^5-4x+2 \in \mathbb{Q}[x]\) is the symmetric group of degree~\(5\).
\item
Let~\(F\) be a free group on a finite set~\(X\) of size~\(n\). Prove that the quotient of~\(F\) by its commutator subgroup is isomorphic to \(\mathbb{Z}^n\).
\item
Let~\(R\) be a ring with~\(1\) and let 
\[
0 \to A \xrightarrow{f} B \xrightarrow{g} C \to 0
\]
 be a short exact sequence of (left)~\(R\)-modules. Show that for any (right)~\(R\)-module~\(M\), the sequence 
\[
M \otimes_R A \xrightarrow{1^M \otimes f} M \otimes_R B \xrightarrow{1^M \otimes g} M \otimes_R C \to 0
\]
 of tensor products, is exact.
\item
Let~\(F\) be a field and \(F[[x]]\) the ring of formal power series in the indeterminate~\(x\).
\begin{examsubparts}
\item[\textnormal{(a)}]
Prove that \(F[[x]]\) is a principal ideal domain.
\item[\textnormal{(b)}]
Determine all of the prime ideals of \(F[[x]]\).
\item[\textnormal{(c)}]
Is \(F[[x]]\) Noetherian? Artinian?
\end{examsubparts}
\item
All parts involve ideals in a commutative ring with identity. Let~\(S\) be a multiplicative set in~\(R\).
\begin{examsubparts}
\item[\textnormal{(a)}]
Prove that if~\(I\) is an ideal of~\(R\) then \(S^{-1}I=\left\{\frac{a}{s}\mid a\in I, s\in S\right\}\) is an ideal in the ring \(S^{-1}R\), the localization of~\(R\) with respect to~\(S\).
\item[\textnormal{(b)}]
Define the radical \(\operatorname{Rad} I\) of an ideal~\(I\) and show that it is an ideal.
\item[\textnormal{(c)}]
Prove that \(S^{-1}(\operatorname{Rad} I)=\operatorname{Rad}(S^{-1}I)\).
\end{examsubparts}
\item
You may not answer this problem just by quoting Nakayama’s Lemma. Let~\(R\) be a commutative local ring with~\(1\) and~\(J\) its unique maximal ideal.
\begin{examsubparts}
\item[\textnormal{(a)}]
Prove that \(1-a\) is a unit for every element~\(a\) of~\(J\),
\item[\textnormal{(b)}]
Let~\(M\) be an~\(R\)-module, with generators \(x_1,x_2,x_3\) and suppose that \(x_2\) and \(x_3\) lie in \(JM\). Show that \(M=Rx_1\).
\end{examsubparts}
\item
Let~\(R\) be a ring with~\(1\). All modules are unitary.
\begin{examsubparts}
\item[\textnormal{(a)}]
Define the concept of a projective~\(R\)-module.
\item[\textnormal{(b)}]
Show that if~\(P\) is a projective~\(R\)-module, then there is an~\(R\)-module~\(K\) such that \(P\oplus K\) is isomorphic to a free module.
\item[\textnormal{(c)}]
Show that for every nonzero element~\(x\) of a free~\(\mathbb{Z}\)-module~\(F\) there is a homomorphism from~\(F\) to a finite group such that the image of~\(x\) is nontrivial.
\item[\textnormal{(d)}]
Prove that~\(\mathbb{Q}\) is not a projective~\(\mathbb{Z}\)-module. (Hint: Show that~\(\mathbb{Q}\) has no nontrivial finite quotients. Then use the earlier parts.)
\end{examsubparts}
\item
Let~\(D\) be a division ring with center~\(K\) and suppose~\(a\) and~\(b\) are elements of~\(D\) which are algebraic over~\(K\) and have the same minimum polynomial. Prove that \(b=dad^{-1}\) for some nonzero element \(d\in D\).
\item
Let~\(\mathcal{C}\) be a category and let \((A_i)_{i\in I}\) be a family of objects of~\(\mathcal{C}\).
\begin{examsubparts}
\item[\textnormal{(a)}]
Define what is a coproduct of \((A_i)_{i\in I}\) in~\(\mathcal{C}\).
\item[\textnormal{(b)}]
Prove that in the category of sets the coproduct of two sets is their disjoint union.
\end{examsubparts}
\end{examproblems}
\end{document}
